\title[Finite failure schemes and quadratic pencils]{Finite failure schemes and the reconstruction of quadratic pencils}
\author{Qian Qi}
\date{September 24, 2026}
\subjclass[2020]{14M12, 14B05, 14D20, 13C40}
\keywords{Quadratic pencil, nonreduced degeneracy locus, intrinsic reconstruction, infinitesimal neighbourhood}
\hypersetup{pdftitle={Finite failure schemes and the reconstruction of quadratic pencils, revision 146},pdfauthor={Qian Qi}}
\begin{abstract}
We reconstruct every complex quadratic pencil from an abstract unmarked
finite-order neighbourhood of its multiplication-failure scheme.  For
embedding dimension \(n\ge3\), the smallest uniform order is
\(d=n^2+2n-4\); all lower orders are independent of the pencil.  A
neighbourhood supported on one Schubert line already suffices.  More
generally, for a pencil subbundle varying over a smooth connected
projective base with an arbitrary source bundle, the unmarked
finite scheme recovers the base, the actual source bundle and the whole
pencil family, up to one constant left transformation.  The proof
recovers intrinsic tensor rulings and the varying coefficient line.
We describe all geometric isomorphisms through their linear data and a
filtered nilpotent automorphism kernel.  Coefficient-support asymmetry
also yields reconstruction from a single unmarked Artin local algebra,
whose tangent-identity automorphism group is explicitly unipotent.  An explicit nonisotrivial
family illustrates reconstruction through spectral collisions.  The
universal finite neighbourhood is compatible with arbitrary complex
base change with its relative socle retained.  For regular pencils the
spectral readout includes all elementary divisors and higher Fitting
incidence schemes.  Framed parameter embeddings, unmarked geometric
inverses and relative constructions are distinguished throughout.
\end{abstract}
\maketitle
